\section{Revealed-Sacrifice Observation Model}
\label{sec:sacrifice_model}

The privacy-compatible observation channel for $\volR$ is
\emph{revealed sacrifice}: agents disclose the value they place on
unbenchmarked capabilities through the benchmarked capabilities
they voluntarily surrender to obtain them.  The framework does not
observe \emph{what} the agent values privately---only \emph{what
they commit to publicly through sacrifice}.

\subsection{Core definitions}

\begin{definition}[Revealed-Sacrifice Event]
\label{def:sacrifice_event}
A revealed-sacrifice event is a tuple $(i, X, Y, t)$ where:
\begin{itemize}
\item $i$ is the agent performing the sacrifice;
\item $X$ is a benchmarked capability (or capability bundle) with
  known $\Delta\volL(X) \geq 0$, surrendered by agent~$i$ at
  time~$t$;
\item $Y = \{Y_1, \ldots, Y_k\}$ is an unbenchmarked capability
  bundle acquired by agent~$i$, satisfying the \emph{genuine-
  exchange conditions}:
  (i)~$X \cap Y = \emptyset$ (the agent does not surrender and
  reacquire the same capability);
  (ii)~$Y \cap (H_i \setminus X) = \emptyset$ (the acquired
  bundle is novel relative to the agent's retained holdings
  $H_i \setminus X$, where $H_i$ is the agent's pre-trade
  holdings);
  (iii)~\emph{post-trade cooperative closure}: no cooperative
  capability in the poset has participants spanning both $Y$
  and $H_i \setminus X$ (so post-trade, no new cross-bundle
  cooperative capability is activated by the combination);
  and
  (iv)~\emph{pre-trade cooperative closure}: no cooperative
  capability in the poset has participants spanning both $X$
  and $H_i \setminus X$ (so pre-trade, no cooperative
  capability straddles the surrendered and retained
  capability sets---only cooperatives contained entirely in
  $X$ or entirely in $H_i \setminus X$ exist);
\item $t$ is the event timestamp.
\end{itemize}
Conditions (i) and (ii) characterize \emph{genuine} trades: the
agent exchanges distinct capabilities and acquires something
novel.  Conditions (iii) and (iv) together ensure that the
exchange-to-bundle cancellation in the revealed-preference
derivation is exact: (iii) rules out post-trade cross-bundle
cooperatives that would inflate $\Ui((H \setminus X) \cup Y)$
beyond $\Ui(H \setminus X) + \Ui(Y)$; (iv) rules out pre-trade
cross-bundle cooperatives that would inflate $\Ui(H)$ beyond
$\Ui(H \setminus X) + \Ui(X)$.  Both closures are needed
symmetrically for the cancellation of
Lemma~\ref{lem:cancellation} to hold.  Trades violating
(i)~or~(ii) are degenerate; trades violating (iii)~or~(iv) carry
revealed-preference information about the marginal contribution
of $Y$ in context rather than $\volR(Y)$ in isolation, but the
theorem's isolated-$\volR(Y)$ interpretation and
Lemma~\ref{lem:cancellation} require all four conditions.
Contextual-quantity salvage is outside this paper's formal
scope and is flagged as a future-work direction in
Section~\ref{sec:discussion}.
The event emits the signal:
\[
  \volRlower(Y) \;\geq\; \Delta\volL(X)
\]
under assumptions S0 (calibration), S1 (free choice), S2
(benchmark grounding), and S4(a) (valuation additivity) stated
below.  S4(a) is needed for the exchange-to-bundle cancellation
(Lemma~\ref{lem:cancellation}).
\end{definition}

\begin{lemma}[Exchange-to-Bundle Cancellation]
\label{lem:cancellation}
Let $H$ be a finite capability holding set with $X \subseteq H$,
and let $Y$ be a finite capability set satisfying the
genuine-exchange conditions (i)--(iv) of
Definition~\ref{def:sacrifice_event} relative to $H$:
(i)~$X \cap Y = \emptyset$;
(ii)~$Y \cap (H \setminus X) = \emptyset$;
(iii)~no cooperative capability in the poset has participants
spanning both $Y$ and $H \setminus X$;
(iv)~no cooperative capability in the poset has participants
spanning both $X$ and $H \setminus X$.
Under S4(a) (global valuation additivity:
$\Ui(S) = \sum_{c \in S} u_i(c)$ for any finite
$S \subseteq P$, with cooperative capabilities included as
explicit nodes in $S$ when their participants are all in $S$):
\[
  \Ui\bigl((H \setminus X) \cup Y\bigr) - \Ui(H)
  \;=\; \Ui(Y) - \Ui(X).
\]
Equivalently,
$\Ui((H \setminus X) \cup Y) \geq \Ui(H)
\iff \Ui(Y) \geq \Ui(X)$.
\end{lemma}

\begin{proof}
By S4(a), $\Ui(A) = \sum_{c \in A} u_i(c)$ for any finite
capability set $A$.  To apply this cleanly we must track
cooperative capabilities carefully: under S4(a), a cooperative
capability $c_{\mathrm{coop}}$ contributes $u_i(c_{\mathrm{coop}})$
to $\Ui(A)$ whenever $c_{\mathrm{coop}} \in A$, which holds iff
all of $c_{\mathrm{coop}}$'s participants are in $A$.

\emph{Pre-trade state $H$.}  $H = (H \setminus X) \sqcup X$.
By (iv), no cooperative capability has participants spanning
$X$ and $H \setminus X$, so every cooperative capability in $H$
has participants entirely in $X$ or entirely in $H \setminus X$.
Hence
\[
  \Ui(H) = \Ui(H \setminus X) + \Ui(X),
\]
where each term includes cooperatives formable within that set
alone.

\emph{Post-trade state $(H \setminus X) \cup Y$.}  By (ii) this
is a disjoint union: $(H \setminus X) \sqcup Y$.  By (iii),
no cooperative capability has participants spanning $Y$ and
$H \setminus X$, so every cooperative in the post-trade state
has participants entirely in $Y$ or entirely in $H \setminus X$.
Hence
\[
  \Ui\bigl((H \setminus X) \cup Y\bigr)
  = \Ui(H \setminus X) + \Ui(Y).
\]

\emph{Cancellation.}
Subtracting, the $\Ui(H \setminus X)$ terms cancel, leaving
$\Ui((H \setminus X) \cup Y) - \Ui(H) = \Ui(Y) - \Ui(X)$.

Condition (i) makes $\Ui(Y)$ and $\Ui(X)$ refer to distinct
capability sets; violating (i) would cause $X \cap Y$ terms to
cancel trivially on both sides, yielding no information.
\end{proof}

\paragraph{Justification (revealed-preference inequality).}
The transfer from the observed trade to the framework-level
bound $\volRlower(Y) \geq \Delta\volL(X)$ proceeds in three
steps.

\begin{enumerate}
\item \textbf{Exchange preference (revealed-preference
  observation).}
  When agent~$i$ trades $X$ for $Y$, revealed-preference
  theory~\citep{samuelson1938note} establishes that the
  post-trade state is weakly preferred to the pre-trade state:
  $\Ui\bigl((H \setminus X) \cup Y\bigr) \geq \Ui(H)$,
  where $H$ is the agent's pre-trade holdings.

\item \textbf{Reduction to bundle-level comparison.}
  By Lemma~\ref{lem:cancellation} under S4(a) and the
  genuine-exchange conditions of
  Definition~\ref{def:sacrifice_event}, the exchange comparison
  reduces to $\Ui(Y) \geq \Ui(X)$.

\item \textbf{Calibration chain.}
  Given $\Ui(Y) \geq \Ui(X)$, the bridge to the framework-level
  bound uses:
  \begin{itemize}
  \item Sacrifice-side grounding (S2):
    $\Ui(X) \geq \Delta\volL(X)$, which holds by construction
    when $\volL$ is the operational target on the benchmarked
    subspace.
  \item Acquisition-side calibration (S0):
    $\Ui(Y) \leq \volR^{[W]}(Y)$, asserting that agents do not
    systematically overvalue unbenchmarked bundles relative to
    their \emph{window-integrated} exercise contribution.  The
    window-integrated form is what the observation channel
    actually bounds: an acquisition at $t_n$ carries no
    instantaneous exercise witness, but the bound is still
    valid on the cumulative window quantity $\volR^{[W]}$
    whenever the agent uses $Y$ during $W$.
  \end{itemize}
  Chaining: $\volR^{[W]}(Y) \geq \Ui(Y) \geq \Ui(X) \geq
  \Delta\volL(X)$.
\end{enumerate}

\begin{remark}[S4(a) is a genuine assumption, not a consequence]
\label{rem:s4a_status}
S4(a) is a global additive-separability assumption on $\Ui$.
It is stronger than the rationality-only premises of Samuelson's
original revealed-preference
theorem~\citep{samuelson1938note}, which establishes preferences
over states rather than additivity over bundles.  In applied
microeconomics~\citep{mascolell1995microeconomic}, bundle-level
inference of the kind the sacrifice channel performs typically
invokes either additive separability or quasilinearity in a
numeraire.  The paper uses a special case implicitly via the
time-sacrifice channel (Definition~\ref{def:time_sacrifice}),
which values time at wage rate $r_i$---quasilinearity in
wage-scaled time, a form of valuation additivity.  S4(a) is the
explicit assumption that makes Lemma~\ref{lem:cancellation} go
through; it is not a consequence of S0--S3.

Complementarities that would break pure additivity in general
consumer theory are absorbed into the poset as cooperative
nodes via M6~\citep[Proposition~1]{lasser2026poset}: cooperative
capabilities are themselves capabilities with their own weights,
so bundle-specific cooperative contributions appear as separate
summands in $\Ui$.  S4(a)'s additivity is compatible with this
structure because cooperative terms are explicit capabilities,
not hidden non-linearities.
\end{remark}

\subsection{Assumptions}
\label{sec:assumptions}

The revealed-sacrifice channel and its aggregate bound
(Section~\ref{sec:lower_bound}) operate under six named
assumptions.  We collect them here for single-point reference.

\begin{description}
\item[S0 (Valuation-bounded-by-Exercise).] For each agent~$i_n$
  with event at $t_n \in W$, the agent's valuation of the
  unbenchmarked bundle $Y_n$ does not exceed $Y_n$'s
  \emph{window-integrated} realized $\volR$ contribution:
  $\Uin{n}(Y_n) \leq \volR^{[W]}(Y_n)$.
  This is the \emph{substantive bridge} between valuation (what
  the agent pays for at $t_n$) and realized exercise (what the
  framework cares about, aggregated over the observation window
  $W$).  It is not a calibration-in-the-narrow-sense (unbiased
  estimation) but a one-sided conservatism claim: agents do not
  systematically pay more than what the capability will actually
  deliver in exercise \emph{within $W$}.  The window-integrated
  form is essential: the sacrifice channel observes acquisition
  at $t_n$, not exercise at any particular time, so S0's
  consequent is about the cumulative exercise witness across the
  window rather than the instantaneous value at $t_n$.  Agents
  who buy capabilities they never use within $W$ violate S0 for
  those events; this matches the intuition that the lower bound
  is about \emph{realized} rather than \emph{stored} capability.

\item[S1 (Free choice).] Each sacrifice event is a voluntary
  trade: agent~$i_n$ chose $Y_n$ over $X_n$ when both were
  available.  The agent had \emph{non-degrading} alternatives:
  retaining $X_n$ was a viable option that would not have left
  the agent \emph{below-sufficiency} on any need dimension.
  For this paper we treat ``below-sufficiency'' as a primitive:
  an event is below-sufficiency iff retaining $X_n$ would leave
  the agent unable to function on at least one need dimension
  without corrective intervention (food, water, shelter,
  medical care, etc.).  S1 fails for events in which retaining
  $X_n$ would put the agent below-sufficiency, and the
  institutional attestation layer
  (Remark~\ref{rem:external_attestation}) is responsible for
  filtering such events from the aggregate.  The full
  poset-construction machinery that identifies below-sufficiency
  structurally---multi-dimensional need bundles, downstream-safe
  sufficiency, polarity-correct benchmarks, the downstream-gated
  weighting that makes need-satisfaction a precondition for
  $\volL$-registration of downstream capability weight---is
  developed in the companion paper~\citep{lasser2026paper8b}.
  This paper does not depend on that machinery for any proof;
  the companion paper operationalizes what this paper
  axiomatizes as a primitive.

\item[S2 (Benchmark grounding).] The benchmarked sacrifice $X_n$
  has a well-defined $\volL$ contribution $\Delta\volL(X_n)$
  computed via~\citep[Proposition~1]{lasser2026poset}, and the
  agent's valuation of $X_n$ is at least as large:
  $\Uin{n}(X_n) \geq \Delta\volL(X_n)$.

\item[S3 (Bundle coherence).] The hedonic decomposition of each
  bundle $Y_n$ into capabilities in $U$ is well-defined (no
  multicollinearity degeneracy).

\item[S4 (Additive separability).] \mbox{}
  \begin{enumerate}
  \item[(a)] \emph{Valuation additivity (global):}
    $\Ui$ is additive over any finite capability set:
    $\Ui(S) = \sum_{c \in S} u_i(c)$ for any finite
    $S \subseteq P$, with $u_i(c) \geq 0$.  In particular this
    applies to sacrifice bundles ($\Ui(X) = \sum_{c \in X}
    u_i(c)$), acquired bundles ($\Ui(Y)$), and agent
    holdings ($\Ui(H)$).  Complementarities are handled at the
    poset level by treating cooperative capabilities as their
    own nodes (M6), which appear as additional capabilities in
    $S$ whose $u_i$ accounts for the cooperative term.
  \item[(b)] \emph{$\volR$ additivity on poset-independent
    capabilities:} The $\volR$ contribution of a set of
    capabilities that are pairwise poset-independent (no
    subsumption relations \emph{and} no cooperative composition
    links) is additive.
  \end{enumerate}

\item[S5 (Decomposition validity).] For events with
  $\volR(Y_n) > 0$, the event-local decomposition
  coefficients are \emph{lower bounds} on the true $\volR$
  share:
  $\alpha_{n,c} \leq \volR(c) / \volR(Y_n)$ for all such $n$
  and $c \in Y_n$.  Events with $\volR(Y_n) = 0$ carry no
  share information and are excluded from Part~B by the
  positivity precondition of the proof.
\end{description}

\begin{remark}[S5 is a conditional-refinement assumption]
\label{rem:s5_heavy}
S5 is the strongest of the six assumptions: it asserts a direct
share inequality on the quantity Part~B of
Theorem~\ref{thm:aggregate_bound} bounds.  Part~A makes no
per-capability attribution claim and does not require S5.
Part~B's per-capability refinement should be understood as a
\emph{conditional result}: under the additional structural
assumption that decomposition coefficients are conservative
share lower bounds, \emph{and} the hypothesis of within-bundle
poset-independence (Definition~\ref{def:bundle_decomp},
Remark~\ref{rem:within_bundle_indep}), the aggregate can be
refined to per-capability bounds.  Within-bundle
poset-independence is what makes S5's share inequality
compatible with the partition-of-unity constraint on the
$\alpha$ coefficients: bundles with internal cooperative
activation produce $\volR(Y_n) > \sum_{c \in Y_n} \volR(c)$,
forcing $\sum_{c \in Y_n} \alpha_{n,c} < 1$ under S5 and
violating Definition~\ref{def:bundle_decomp}.  Operationally, S5 holds when
decomposition coefficients are set via conservative heuristics
(e.g., equal-weight default $\alpha_{n,c} = 1/|Y_n|$ when $Y_n$
is a small homogeneous bundle, known-bias corrections from prior
hedonic~\citep{rosen1974hedonic} regressions) rather than via
fitted models that could overstate individual capability
contributions.
\end{remark}

Part~A of Theorem~\ref{thm:aggregate_bound} requires
S0--S2 and S4(a): S4(a) is the valuation-additivity assumption
used in the exchange-to-bundle reduction of the revealed-
preference inequality (Step~2 of the justification above).
Part~B additionally requires S3 (bundle coherence), S4(b)
($\volR$ additivity on poset-independent capabilities), S5
(decomposition validity), and within-bundle poset-independence
of each admissible $Y_n$ (Definition~\ref{def:bundle_decomp}).
The need-sufficiency architecture of the companion
paper~\citep{lasser2026paper8b} characterizes when S1 fails
structurally and provides the filter that keeps the sacrifice
channel clean; this paper treats the filter as an external
precondition supplied by Remark~\ref{rem:external_attestation}.

\begin{remark}[S0 is the substantive bridge, not a calibration
detail]
\label{rem:s0_substantive}
S0 does the conceptual work of the theorem.  Revealed-preference
reasoning on its own establishes only that traded bundles are
valued at least as highly as surrendered bundles; it does not
establish that what the agent values corresponds to what the
capability \emph{delivers in exercise}.  S0 asserts precisely
that correspondence (in the one-sided direction needed for a
lower bound): subjective valuation is upper-bounded by realized
exercise contribution.  This is a non-trivial behavioural claim.
Agents who systematically overpay for status, novelty, or
anticipated-but-unrealized utility violate S0 in realistic
markets; the theorem's empirical validity depends on the
prevalence of such violations.

S0 does \emph{not} require that agent valuations perfectly track
$\volR$---only that they do not \emph{overstate} it.  Agents may
undervalue unbenchmarked bundles (sacrifice little for something
very valuable), in which case the lower bound is conservative
but still valid.  S0 fails only in the overpayment direction.
Weakening S0 to expected/distributional conservatism
(Open Question~\ref{oq:calibration}) is a priority direction for
follow-up work.
\end{remark}

\begin{remark}[Common-scale convention for the calibration chain]
\label{rem:units}
The chain of inequalities $\volR^{[W]}(Y) \geq \Ui(Y) \geq \Ui(X)
\geq \Delta\volL(X)$ presupposes a common scale on which the four
quantities are comparable.  We adopt the \emph{poset-measure
scale}: $\volL$ is the framework's measure-theoretic primitive
(\citep[Proposition~1]{lasser2026poset}), and $\volR$ inherits its
unit from $\volL$ on the exercised sub-poset
(Definition~\ref{def:volR}).  S0 is the assumption that each
agent's subjective valuation $\Ui$ is reported in this scale (i.e.,
$\Ui(Y)$ has the same units as $\volR^{[W]}(Y)$); S2 is the
analogous assumption on the benchmarked side ($\Ui(X)$ in the
same units as $\Delta\volL(X)$).  Money price $p$ enters through
the money-sacrifice channel
(Definition~\ref{def:money_sacrifice}) only via
$\Delta\volL(X(p))$, which converts price to poset-measure units
through the benchmark calibration of~\citep[Definition~2]{lasser2026poset};
analogously, the time-sacrifice channel
(Definition~\ref{def:time_sacrifice}) converts the wage-valued
quantity $r_i \cdot h$ to poset-measure units through
$\Delta\volL(r_i \cdot h)$.  Currency and wage-time are not
themselves on the poset-measure scale; they are inputs to the
benchmark function whose output is.  The theorem's inequalities
hold in this single common scale; cross-currency or cross-wage
comparability is a property of the benchmark function, not of the
theorem.
\end{remark}

\begin{definition}[Aggregate Trade Window]
\label{def:trade_window}
An aggregate trade window $[t_0, t_0 + T]$ is an observation
period over which revealed-sacrifice events are collected.  The
window parameters are:
\begin{itemize}
\item $T$: window duration;
\item $N$: number of observed sacrifice events in the window;
\item $\{(i_n, X_n, Y_n, t_n)\}_{n=1}^N$: the event sequence.
\end{itemize}
The aggregate trade window plays a role analogous to the
observation windows used elsewhere in the
sequence~\citep{lasser2026scm}: it defines the timescale
over which the framework accumulates sacrifice evidence.
\end{definition}

\subsection{Privacy properties}

The channel's privacy properties split cleanly into two
categories: \emph{formal} properties that follow directly from
the observation-model construction (stated below as
Theorem~\ref{thm:privacy}), and \emph{institutional}
preconditions that require deployment-level infrastructure to
hold operationally (stated as
Proposition~\ref{prop:privacy_institutional}).  We state each
category separately so that the theorem-level content covers
only mathematical consequences of the construction, while
institutional dependencies are flagged as preconditions rather
than proved facts.

\begin{theorem}[Formal Privacy Properties of Revealed-Sacrifice
Observation]
\label{thm:privacy}
The revealed-sacrifice channel satisfies the following three
formal privacy properties, each following from the observation
tuple (Definition~\ref{def:sacrifice_event}) together with the
commitment protocol
of~\citep[Definition~6]{lasser2026exo}:

\begin{enumerate}
\item[\textbf{(P1)}] \textbf{No interior access.}  The framework
  never observes what the agent values privately---only what they
  commit to publicly through sacrifice.  The observation channel
  records the tuple $(i, X, Y, t)$, not the agent's internal
  valuation function $\Ui$.

\item[\textbf{(P3)}] \textbf{Discretization as a feature.}  Trade
  events are sparse by construction.  The channel's emission rate
  is bounded by trade frequency, not by observation intensity:
  each sacrifice event produces one output tuple, and no
  continuous stream is emitted between events.  (This is an
  emission-rate bound, not a Shannon-information-rate bound; the
  information content per event is finite but non-trivial.)

\item[\textbf{(P4)}] \textbf{Commitment-layer composability.}
  The commitment protocol
  of~\citep[Definition~6]{lasser2026exo} extends to
  revealed-sacrifice events: a sacrifice event can be committed as
  a proof-of-trade exposing only the $\volR$-category and
  $\Delta\volL$-magnitude, not counterparty, price, or specific
  goods (Section~\ref{sec:commitment}).
\end{enumerate}
\end{theorem}

\begin{proof}
Each property is established from the structure of the observation
channel (Definition~\ref{def:sacrifice_event}) and the aggregate
bound (Theorem~\ref{thm:aggregate_bound}).

\emph{(P1) No interior access.}
The channel's output is, by Definition~\ref{def:sacrifice_event},
the tuple $(i, X, Y, t)$.  The agent's internal valuation $\Ui$
appears only as a bridge quantity in the calibration chain
(S0 + S1 + S2 + S4(a)); it is never recorded, transmitted, or
stored by the framework.  The aggregate bound
(Theorem~\ref{thm:aggregate_bound}) depends only on
$\{(\Delta\volL(X_n), Y_n)\}$, so the framework's behaviour is
identical under any two agent-valuation functions producing the
same trade sequence.  No interior access is required or obtained.

\emph{(P3) Discretization as a feature.}
The channel emits at most one output tuple per sacrifice event
and produces no continuous stream otherwise.  Over a window of
duration $T$ the total number of channel emissions is exactly
the number of observed events, so the channel's \emph{emission
rate} is $|\mathrm{events}|/T$.  Unlike continuous observation
(which scales with sampling rate), the emission rate is bounded
above by the agents' trade frequency---a behavioural rather than
observational parameter---so increasing observation intensity
does not increase the emission rate.  (Each emission carries
finite information about the tuple $(i, X, Y, t)$; this is an
emission-rate claim, not a Shannon-information-rate claim.)

\emph{(P4) Commitment-layer composability.}
Each sacrifice event admits a commitment
$\Gamma = \Commit\bigl((\Delta\volL(X), \catY), r\bigr)$
using the protocol of~\citep[Definition~6]{lasser2026exo},
accompanied by a zero-knowledge
proof~\citep[Definition~7]{lasser2026exo} that $\Gamma$ binds a
valid revealed-sacrifice event: the event-local objectively
verifiable facts---non-negativity of $\Delta\volL(X_n)$, bundle
containment $Y_n \subseteq C_{\catYn{n}}$
(Definition~\ref{def:committed_event}(d)), bundle coherence S3,
and a verifiable link to an underlying trade record.  This is
the event-level certification; S4(b) is not an event-level
property but a batch-level structural fact about the category
partition, established via pairwise-distinct category labels in
Proposition~\ref{prop:zk_aggregation} and inherited from the
poset-independence of distinct partition components
(Definition~\ref{def:category_partition}(i)).  S0, S1, S4(a),
and S5 are \emph{not} ZK-verifiable
(Remark~\ref{rem:zk_limits}).
Proposition~\ref{prop:commit_preserves} establishes that the
Part~A aggregate bound is computable from committed events with
\emph{pairwise-distinct} category labels, together with the
category-partition structure; batches with colliding labels
require the operational responses of
Remark~\ref{rem:label_collision} (sub-batch partitioning, finer
public partition, or richer commitments).  The committed form
reveals only the sacrifice magnitude and bundle category; agent
identity, specific $X$ and $Y$, and counterparty remain hidden
by the hiding property of
$\Commit$~\citep{pedersen1991noninteractive,lasser2026exo}.
\end{proof}

\begin{proposition}[Institutional Preconditions for
Revealed-Sacrifice Observation]
\label{prop:privacy_institutional}
Two additional privacy properties hold operationally when the
deployment environment supplies the corresponding institutional
infrastructure.  These are not formal consequences of the
channel's construction; they are preconditions the framework
relies on but does not itself construct.

\begin{enumerate}
\item[\textbf{(P2)}] \textbf{Consent via participation
  (institutional).}
  The aggregate lower bound sums per-capability max-attribution
  quantities over observed events
  (Theorem~\ref{thm:aggregate_bound}, Part~B), so an agent with
  zero observed sacrifice events contributes zero to the
  aggregate.  This zero-contribution property is formal.  The
  \emph{operational} interpretation as ``consent'' additionally
  requires that participation in markets and labour exchanges be
  legally and socially opt-in; in settings where participation
  is coerced, the formal zero-contribution property holds but
  the consent interpretation does not.

\item[\textbf{(P5)}] \textbf{Third-party observability
  (institutional).}
  \emph{Money-sacrifice events} are observable from
  public-facing ledger data (market exchanges, payment rails,
  public commitments) with no agent participation required,
  \emph{assuming} such public-ledger infrastructure exists.  The
  price~$p$ is directly readable from public ledger data; the
  sacrifice magnitude $\Delta\volL(X)$ is then computed from~$p$
  via the benchmark function on the surrendered purchasing
  power (Definition~\ref{def:money_sacrifice}).  ``Directly
  readable'' refers to~$p$; $\Delta\volL$ is a function of~$p$
  through a public benchmark, not a separately-disclosed
  quantity.

  \emph{Time-sacrifice events} are observable from public-ledger
  data only under an additional precondition: the agent's
  outside-option wage rate $r_i$ (Definition~\ref{def:time_sacrifice})
  must be computable from public inputs---e.g., a public wage
  schedule, occupational census with public imputation rules, or
  the agent's own published rate.  Where $r_i$ is private, the
  hours $h$ may be observable but the sacrifice magnitude
  $\Delta\volL(r_i \cdot h)$ is not third-party computable;
  P5 then holds for the money channel and for
  public-wage-imputation time events, but not for fully private
  time-sacrifice valuations.

  Where public-ledger infrastructure does not exist at all, P5
  fails operationally on both channels even though the
  construction is unchanged.  This contrasts with controlled
  relaxation~\citep{lasser2026wamura} (framework must impose a
  test) and with commitment-protocol
  verification~\citep{lasser2026exo} (agent must generate a
  proof).  P5 does not extend to private time-sacrifice
  activities (parenting, contemplation), which require voluntary
  disclosure regardless of ledger infrastructure and regardless
  of wage imputation.
\end{enumerate}
\end{proposition}

\begin{remark}[Why the split]
\label{rem:privacy_mixed_status}
Theorem~\ref{thm:privacy} restricts itself to properties that are
mathematical consequences of the channel's construction.
Proposition~\ref{prop:privacy_institutional} catalogues the
deployment conditions under which P2 and P5 hold operationally.
Treating P2 and P5 as institutional preconditions rather than as
theorem consequences is the honest framing: they are load-bearing
for how the channel behaves in practice, but they depend on facts
about markets, legal systems, and public infrastructure that lie
outside the paper's formal scope.
\end{remark}
